
\documentclass[tikz,border=8pt]{standalone}
\usepackage[UTF8,fontset=fandol]{ctex}
\usepackage{amsmath,amssymb}
\usetikzlibrary{arrows.meta,calc,positioning}
\definecolor{ink}{HTML}{172033}
\definecolor{muted}{HTML}{637083}
\definecolor{blue}{HTML}{4F86C6}
\definecolor{green}{HTML}{61A66A}
\definecolor{red}{HTML}{D55E5E}
\definecolor{amber}{HTML}{D59B3D}
\definecolor{paperblue}{HTML}{F2F7FD}
\definecolor{papergreen}{HTML}{F3F9F1}
\definecolor{paperamber}{HTML}{FFF8E8}
\tikzset{
  panel/.style={anchor=north west,minimum width=9.15cm,minimum height=6.05cm,
    text width=8.45cm,align=left,inner sep=10pt,draw=black!12,rounded corners=5pt,fill=white},
  title/.style={font=\bfseries\large,text=ink,anchor=west},
  body/.style={font=\small,text=muted,align=left,text width=7.9cm},
  formula/.style={draw=blue!35,fill=paperblue,rounded corners=4pt,inner sep=8pt,
    text=ink,align=center,text width=7.7cm},
  insight/.style={draw=green!40,fill=papergreen,rounded corners=4pt,inner sep=7pt,
    text=ink,align=left,text width=7.7cm,font=\small},
  arrow/.style={-{Latex[length=2.3mm]},line width=1pt,draw=blue!70},
  chip/.style={rounded corners=3pt,draw=black!14,fill=black!2,inner xsep=7pt,
    inner ysep=5pt,font=\small,text=ink},
}
\newcommand{\Panel}[5]{%
  \node[panel] (#1) at (#2,#3) {};
  \node[title] at ($(#1.north west)+(0.35,-0.42)$) {#4};
  \node[body,anchor=north west] at ($(#1.north west)+(0.35,-1.05)$) {#5};
}
\newcommand{\Summary}[1]{%
  \node[draw=amber!35,fill=paperamber,rounded corners=6pt,minimum width=28.2cm,
    minimum height=1.55cm,text width=27.2cm,align=center,font=\small,text=ink]
    at (14.125,-13.35) {\textbf{一图总结}\quad #1};
}
\newcommand{\Identity}{%
  \node[font=\scriptsize,text=black!38,anchor=east] at (28.15,-14.35)
    {cs229-storyboard@五道口纳什};
}
\newcommand{\Formula}[1]{%
  \par\vspace{5pt}\begingroup\setlength{\fboxsep}{8pt}%
  \colorbox{paperblue}{\parbox{7.45cm}{\centering\color{ink}#1}}%
  \endgroup\par\vspace{6pt}}
\newcommand{\Insight}[1]{%
  \par\vspace{5pt}\begingroup\setlength{\fboxsep}{7pt}%
  \colorbox{papergreen}{\parbox{7.55cm}{\color{ink}#1}}%
  \endgroup\par\vspace{6pt}}
\newcommand{\Chip}[1]{%
  \begingroup\setlength{\fboxsep}{5pt}\colorbox{black!4}{\strut\color{ink}#1}\endgroup}

\begin{document}
\begin{tikzpicture}[x=1cm,y=1cm]
\fill[white] (-0.35,0.35) rectangle (28.6,-14.55);
\Panel{p1}{0}{0}{1. 隐变量让似然难算}{观测到 $x$，却看不到 $z$：\Formula{$\ell(\theta)=\sum_i\log\sum_z p(x_i,z;\theta)$}“对数外还有求和”阻碍直接优化。}
\Panel{p2}{9.55}{0}{2. 引入辅助分布}{为每个样本引入 $q_i(z)$，用 Jensen 不等式构造下界：\Formula{$\ell(\theta)\ge\sum_{i,z}q_i(z)\log\dfrac{p(x_i,z;\theta)}{q_i(z)}$}}
\Panel{p3}{19.1}{0}{3. E 步更新责任度}{固定旧参数，令下界贴住真实似然：\Formula{$q_i(z)=p(z\mid x_i;\theta^{old})$}它回答“每个隐状态应负责多少”。}
\Panel{p4}{0}{-6.45}{4. M 步更新模型}{固定责任度，最大化完整数据的期望对数似然：\Formula{$\theta^{new}=\arg\max_\theta\sum_{i,z}q_i(z)\log p(x_i,z;\theta)$}}
\Panel{p5}{9.55}{-6.45}{5. 为什么不会变差}{E 步让下界与当前似然相等，M 步提高下界：\Insight{$\ell(\theta^{new})\ge\mathcal L(q,\theta^{new})\ge\mathcal L(q,\theta^{old})=\ell(\theta^{old})$。}}
\Panel{p6}{19.1}{-6.45}{6. 以高斯混合为例}{\Chip{E：样本属于各簇的概率}\quad\Chip{M：加权均值与协方差}多次初始化以降低落入差的局部最优的风险。}
\Summary{困难似然 $\rightarrow$ Jensen 下界 $\rightarrow$ E 步估计隐状态 $\rightarrow$ M 步更新参数 $\rightarrow$ 单调改进。}
\Identity
\end{tikzpicture}
\end{document}
